\subsection{Casos de Uso del Administrador de TI}
\label{subsec:casos_uso_admin_ti}

La figura \ref{fig:casos_uso_admin_ti} detalla las funciones del Administrador de TI, responsable del mantenimiento técnico, la gestión de catálogos del sistema y la supervisión de la infraestructura de la plataforma.

% Ajustes globales para esta sección para ahorrar espacio
\begingroup
\footnotesize % Fuente reducida para todas las tablas
\setlength{\parskip}{0pt} % Elimina espacio entre párrafos
\setlength{\belowcaptionskip}{-10pt} % Reduce espacio después del título de la tabla
\renewcommand{\arraystretch}{1.1} % Ajusta el alto de las filas para que no se vea amontonado

% ================= CU70 =================
\begin{table}[ht!]
    \centering
    \caption{Especificación de Caso de Uso CU70: Consultar panel de TI}
    \label{tab:cu70}
    \begin{tabularx}{\textwidth}{|l|p{2cm}|l|X|}
        \hline
        \textbf{ID} & CU70 & \textbf{Nombre} & Consultar panel de TI \\ \hline
        \textbf{Actor} & \multicolumn{3}{X|}{Administrador de TI} \\ \hline
        \textbf{Descripción} & \multicolumn{3}{X|}{Interfaz principal con métricas de salud del sistema, logs y configuración.} \\ \hline
        \textbf{Disparador} & \multicolumn{3}{X|}{El usuario accede a la ruta /admin-ti/dashboard.} \\ \hline
        \textbf{Precondiciones} & \multicolumn{3}{X|}{Sesión activa con privilegios de superusuario.} \\ \hline
        \textbf{Flujo} & \multicolumn{3}{X|}{1. Carga estatus de microservicios. 2. Despliega opciones de gestión.} \\ \hline
        \textbf{Error} & \multicolumn{3}{X|}{Error de comunicación con el servicio de monitoreo.} \\ \hline
        \textbf{Postcondición} & \multicolumn{3}{X|}{Panel técnico visualizado.} \\ \hline
    \end{tabularx}
\end{table}

\vspace{-0.4cm} % Espacio negativo para juntar las tablas

% ================= CU71 =================
\begin{table}[ht!]
    \centering
    \caption{Especificación de Caso de Uso CU71: Gestionar catálogos}
    \label{tab:cu71}
    \begin{tabularx}{\textwidth}{|l|p{2cm}|l|X|}
        \hline
        \textbf{ID} & CU71 & \textbf{Nombre} & Gestionar catálogos del sistema \\ \hline
        \textbf{Actor} & \multicolumn{3}{X|}{Administrador de TI} \\ \hline
        \textbf{Descripción} & \multicolumn{3}{X|}{Control sobre las tablas maestras que alimentan los formularios.} \\ \hline
        \textbf{Disparador} & \multicolumn{3}{X|}{El usuario selecciona "Catálogos" en el menú.} \\ \hline
        \textbf{Precondiciones} & \multicolumn{3}{X|}{N/A.} \\ \hline
        \textbf{Flujo} & \multicolumn{3}{X|}{1. Lista catálogos (Escuelas, Cargos). 2. Usuario selecciona uno.} \\ \hline
        \textbf{Error} & \multicolumn{3}{X|}{Error de permisos en BD.} \\ \hline
        \textbf{Postcondición} & \multicolumn{3}{X|}{Módulos de catálogo listos para operar.} \\ \hline
    \end{tabularx}
\end{table}

\vspace{-0.4cm}

% ================= CU72 =================
\begin{table}[ht!]
    \centering
    \caption{Especificación de Caso de Uso CU72: Catálogo de escuelas}
    \label{tab:cu72}
    \begin{tabularx}{\textwidth}{|l|p{2cm}|l|X|}
        \hline
        \textbf{ID} & CU72 & \textbf{Nombre} & Gestionar catálogo de escuelas/facultades \\ \hline
        \textbf{Actor} & \multicolumn{3}{X|}{Administrador de TI} \\ \hline
        \textbf{Descripción} & \multicolumn{3}{X|}{Alta, baja o modificación de las entidades académicas.} \\ \hline
        \textbf{Disparador} & \multicolumn{3}{X|}{Selección del catálogo de escuelas en CU71.} \\ \hline
        \textbf{Precondiciones} & \multicolumn{3}{X|}{N/A.} \\ \hline
        \textbf{Flujo} & \multicolumn{3}{X|}{1. Muestra lista de escuelas. 2. Operaciones CRUD. 3. Actualización.} \\ \hline
        \textbf{Error} & \multicolumn{3}{X|}{Duplicidad de claves de facultad.} \\ \hline
        \textbf{Postcondición} & \multicolumn{3}{X|}{Base de datos de escuelas actualizada.} \\ \hline
    \end{tabularx}
\end{table}

\vspace{-0.4cm}

% ================= CU75 =================
\begin{table}[ht!]
    \centering
    \caption{Especificación de Caso de Uso CU75: Catálogo de cargos}
    \label{tab:cu75}
    \begin{tabularx}{\textwidth}{|l|p{2cm}|l|X|}
        \hline
        \textbf{ID} & CU75 & \textbf{Nombre} & Gestionar catálogo de cargos \\ \hline
        \textbf{Actor} & \multicolumn{3}{X|}{Administrador de TI} \\ \hline
        \textbf{Descripción} & \multicolumn{3}{X|}{Actualización de los puestos de los presuntos responsables.} \\ \hline
        \textbf{Disparador} & \multicolumn{3}{X|}{Selección del catálogo de cargos en CU71.} \\ \hline
        \textbf{Precondiciones} & \multicolumn{3}{X|}{N/A.} \\ \hline
        \textbf{Flujo} & \multicolumn{3}{X|}{1. Presenta lista de cargos. 2. Agrega o edita. 3. Sincroniza.} \\ \hline
        \textbf{Error} & \multicolumn{3}{X|}{Error al eliminar cargo vinculado a quejas activas.} \\ \hline
        \textbf{Postcondición} & \multicolumn{3}{X|}{Cargos actualizados para formularios.} \\ \hline
    \end{tabularx}
\end{table}

\endgroup
\clearpage